% !TEX root = ../research_vs.tex

\begin{center}
\Large{\textbf{Internet Appendix for ``Does Peer-Reviewed Research Help Predict Stock Returns''}}
\end{center}

% results that are not needed to justify main results go here
% this includes additional color 
\section{Additional Results on Data-Mined Predictability}\label{sec:ia-dm}
% EXHIBIT-NUMBER: Table IA.1
\begin{table}[!h]
  \caption{Out-of-Sample Returns from Mining Accounting Data: 2004-2020}
  \label{tab:dm-sum-2004}
  
  \begin{singlespace}
  \noindent 
  We sort 29,000 accounting ratios each June into 5 bins based on past 30-year long-short returns (in-sample) and compute the mean return over the next year within each bin (out-of-sample). Statistics are calculated by strategy, then averaged within bins, then averaged across sorting years. Decay is the percentage decrease in mean return out-of-sample relative to in-sample. We omit decay for bin 4 because the mean return in-sample is negligible. Out-of-sample returns are calculated using only data from 2004-2020.
  \end{singlespace}
  \begin{centering}
  \vspace{-3ex}
  \par\end{centering}
  \centering{}\setlength{\tabcolsep}{1.0ex} \small
  \begin{center}
  \begin{tabular}{crrrrrrrrrrrr} \toprule 
  % \multicolumn{13}{c}{Panel (b): Out-of-Sample Returns 2004-2020} \\ 
  In- &   & \multicolumn{5}{c}{Equal-Weighted Long-Short Deciles} &   & \multicolumn{5}{c}{Value-Weighted Long-Short Deciles} \\ \cmidrule{3-7}\cmidrule{9-13}Sample &   & \multicolumn{2}{c}{Past 30 Years (IS)} &   & \multicolumn{2}{c}{Next Year (OOS)} &   & \multicolumn{2}{c}{Past 30 Years (IS)} &   & \multicolumn{2}{c}{Next Year (OOS)} \\ \cmidrule{3-4}\cmidrule{6-7}\cmidrule{9-10}\cmidrule{12-13}Bin &   & \multicolumn{1}{c}{Return} & \multicolumn{1}{c}{\multirow{2}[2]{*}{t-stat}} &   & \multicolumn{1}{c}{Return} & \multicolumn{1}{c}{Decay} &   & \multicolumn{1}{c}{Return} & \multicolumn{1}{c}{\multirow{2}[2]{*}{t-stat}} &   & \multicolumn{1}{c}{Return} & \multicolumn{1}{c}{Decay} \\   &   & \multicolumn{1}{c}{(bps pm)} &   &   & \multicolumn{1}{c}{(bps pm)} & \multicolumn{1}{c}{(\%)} &   & \multicolumn{1}{c}{(bps pm)} &   &   & \multicolumn{1}{c}{(bps pm)} & \multicolumn{1}{c}{(\%)} \\ \cmidrule{1-1}\cmidrule{3-4}\cmidrule{6-7}\cmidrule{9-10}\cmidrule{12-13}
  
  % latex table generated in R 4.2.3 by xtable 1.8-4 package
% Mon Mar  3 17:09:07 2025
 1 &  & -59.2 & -3.99 &  & -24.9 & 57.9 &  & -37.3 & -1.88 &  & -4.2 & 88.7 \\ 
  2 &  & -28.1 & -2.29 &  & -9.6 & 65.8 &  & -14.6 & -0.91 &  & -1.1 & 92.5 \\ 
  3 &  & -11.7 & -1.01 &  & 0.1 & 100.9 &  & -4.2 & -0.28 &  & -2.6 & 38.7 \\ 
  4 &  & 1.8 & 0.14 &  & 6.7 &  &  & 5.5 & 0.36 &  & -3.7 &  \\ 
  5 &  & 23.9 & 1.48 &  & 16.3 & 31.8 &  & 25.8 & 1.31 &  & 0.6 & 97.8 \\ 
  
  
  \bottomrule \end{tabular}%   
  \end{center} 
\end{table}


% EXHIBIT-NUMBER: Table IA.2
\begin{table}[!h]
  \caption{Pairwise Correlations of Data-Mined Predictors}
  \label{tab:dm-cor}
  
  \begin{singlespace}
  \noindent Data-mined predictors are represented by strategies with t-statistics greater than 2.0 in at least 10\% of the in-sample periods from Table \ref{tab:dm-oos}.  The table reports percentiles of Pearson correlation coefficients computed over pairwise-complete return observations. 
  \end{singlespace}
  \begin{centering}
  \vspace{-2ex}
  \par\end{centering}
  \centering{}\setlength{\tabcolsep}{1.7ex} \small
  \begin{center}
  % ==== begin paste
  \begin{tabular}{clccccccccccc}
  \toprule
   \multicolumn{11}{c}{Panel (a): Pairwise correlations} \\
  \midrule
  Percentile & & 1 & 5 & 10 & 25 & 50 & 75 & 90 & 95 & 99 \\
  \midrule
  % latex table generated in R 4.5.3 by xtable 1.8-4 package
% Thu Aug 20 10:14:59 2026
 Equal-Weighted &  & -0.41 & -0.23 & -0.15 & -0.04 & 0.06 & 0.18 & 0.32 & 0.41 & 0.62 \\ 
  Value-Weighted &  & -0.34 & -0.20 & -0.14 & -0.06 & 0.02 & 0.11 & 0.22 & 0.30 & 0.58 \\ 
  
  \bottomrule
  \end{tabular}
  \end{center}
\end{table}

% =====================================================
\clearpage
\newpage
\section{Data-Mined Themes in Other Samples}\label{sec:ia-themes}

%%
% EXHIBIT-NUMBER: Table IA.3
\begin{table}[!h]
\caption{Themes from Mining Accounting Ratios in 1990}\label{tab:dm-theme1990}
\small

Table reports the 20 accounting ratio numerator and stock weight (equal- or value-) combinations with the largest mean t-stats using returns in the years 1963-1990 (IS). `ew' is equal-weight, `vw' is value-weight. Strategies are signed to have positive mean returns IS. `Pct Short' is the share of strategies that short the ratio.  `t-stat' and  `Mean Return' are averages across the 65 possible denominators. 
\vspace{2ex}

\centering

\setlength{\tabcolsep}{1.0ex}

\begin{tabular}{lrrrlrr}
\toprule
 & \multicolumn{3}{c}{1963-1990 (IS)} &  & \multicolumn{1}{c}{1991-2004} & \multicolumn{1}{c}{1991-2022} \\ \cmidrule{2-4} \cmidrule{6-7}
Numerator (Stock Weight) & \multicolumn{1}{c}{Pct} & \multirow{2}{*}{t-stat} & \multicolumn{1}{c}{Mean} &  & \multicolumn{2}{c}{Mean Return} \\  
 & \multicolumn{1}{c}{Short} &  & \multicolumn{1}{c}{Return} &  & \multicolumn{2}{c}{OOS / IS} \\ 
\midrule
$\Delta$Capital surplus (ew) & 100 & 5.8 & 0.67 &  & 1.04 & 0.94\\
$\Delta$Common stock (ew) & 100 & 5.8 & 0.69 &  & 0.80 & 0.55\\
$\Delta$Liabilities (ew) & 100 & 5.7 & 0.74 &  & 0.87 & 0.56\\
$\Delta$Inventories (ew) & 100 & 5.4 & 0.65 &  & 1.44 & 0.79\\
$\Delta$Current liabilities (ew) & 100 & 5.4 & 0.60 &  & 1.04 & 0.56\\
$\Delta$Debt in current liab (ew) & 100 & 5.2 & 0.48 &  & 0.30 & 0.31\\
Stock issuance (ew) & 100 & 5.2 & 0.89 &  & 1.03 & 0.80\\
$\Delta$Long-term debt (ew) & 100 & 5.1 & 0.53 &  & 1.31 & 0.75\\
$\Delta$Notes payable st (ew) & 100 & 5.1 & 0.46 &  & 0.17 & 0.25\\
$\Delta$Interest expense (ew) & 100 & 5.1 & 0.58 &  & 1.01 & 0.80\\
$\Delta$PPE net (ew) & 100 & 4.8 & 0.73 &  & 1.41 & 0.75\\
$\Delta$PPE gross (ew) & 100 & 4.7 & 0.73 &  & 1.15 & 0.61\\
Retained earnings restatement (ew) & 100 & 4.6 & 0.54 &  & 1.38 & 0.70\\
$\Delta$Assets (ew) & 100 & 4.5 & 0.73 &  & 1.63 & 0.94\\
Stock repurchases (ew) & 0 & 4.4 & 0.38 &  & 0.27 & 0.63\\
$\Delta$Convertible debt and stock (ew) & 100 & 4.1 & 0.42 &  & 1.47 & 1.18\\
$\Delta$Capital surplus (vw) & 100 & 4.0 & 0.57 &  & 0.72 & 0.64\\
$\Delta$Cost of goods sold (ew) & 100 & 3.9 & 0.49 &  & 1.41 & 0.84\\
Long-term debt issuance (ew) & 88 & 3.9 & 0.48 &  & 1.30 & 0.71\\
$\Delta$Invested capital (ew) & 100 & 3.9 & 0.63 &  & 2.16 & 1.20\\
\bottomrule
\end{tabular}


\end{table}


%%
% EXHIBIT-NUMBER: Table IA.4
\begin{table}
\caption{Themes from Mining Accounting Ratios in 2000}\label{tab:dm-theme2000}
\small

Table reports the 20 accounting ratio numerator and stock weight (equal- or value-) combinations with the largest mean t-stats using returns in the years 1963-2000 (IS).  `ew' is equal-weight, `vw' is value-weight. Strategies are signed to have positive mean returns IS. `Pct Short' is the share of strategies that short the ratio.  `t-stat' and  `Mean Return' are averages across the 65 possible denominators. 
\vspace{2ex}

\centering

\setlength{\tabcolsep}{1.0ex}

\begin{tabular}{lrrrlrr}
\toprule
 & \multicolumn{3}{c}{1963-2000 (IS)} &  & \multicolumn{1}{c}{2001-2004} & \multicolumn{1}{c}{2001-2022} \\ \cmidrule{2-4} \cmidrule{6-7}
Numerator (Stock Weight) & \multicolumn{1}{c}{Pct} & \multirow{2}{*}{t-stat} & \multicolumn{1}{c}{Mean} &  & \multicolumn{2}{c}{Mean Return} \\  
 & \multicolumn{1}{c}{Short} &  & \multicolumn{1}{c}{Return} &  & \multicolumn{2}{c}{OOS / IS} \\ 
\midrule
$\Delta$Inventories (ew) & 100 & 6.9 & 0.77 &  & 0.72 & 0.33\\
$\Delta$Long-term debt (ew) & 100 & 6.4 & 0.60 &  & 0.81 & 0.37\\
$\Delta$Common stock (ew) & 100 & 6.3 & 0.66 &  & 0.81 & 0.46\\
$\Delta$PPE net (ew) & 100 & 6.3 & 0.82 &  & 1.10 & 0.37\\
$\Delta$Current liabilities (ew) & 100 & 6.1 & 0.61 &  & 0.94 & 0.33\\
$\Delta$Interest expense (ew) & 100 & 6.1 & 0.61 &  & 0.45 & 0.58\\
$\Delta$Liabilities (ew) & 100 & 6.0 & 0.71 &  & 0.87 & 0.44\\
$\Delta$PPE gross (ew) & 100 & 5.9 & 0.78 &  & 0.87 & 0.30\\
$\Delta$Debt subordinated convertible (ew) & 100 & 5.4 & 0.71 &  & 1.15 & 0.62\\
$\Delta$Debt convertible (ew) & 100 & 5.4 & 0.61 &  & 1.72 & 0.74\\
Retained earnings restatement (ew) & 100 & 5.4 & 0.61 &  & 1.07 & 0.29\\
$\Delta$Invested capital (ew) & 100 & 5.3 & 0.83 &  & 1.55 & 0.56\\
Merger sales contrib (ew) & 100 & 5.2 & 0.53 &  & 0.93 & 0.51\\
$\Delta$Assets (ew) & 100 & 5.2 & 0.86 &  & 1.33 & 0.53\\
$\Delta$Capital surplus (ew) & 100 & 5.2 & 0.69 &  & 0.86 & 0.85\\
$\Delta$Capital expenditure (ew) & 100 & 5.2 & 0.53 &  & 1.67 & 0.64\\
$\Delta$Cost of goods sold (ew) & 100 & 5.1 & 0.58 &  & 0.66 & 0.40\\
$\Delta$Num employees (ew) & 100 & 5.0 & 0.59 &  & 1.42 & 0.52\\
$\Delta$Intangible assets (ew) & 100 & 5.0 & 0.49 &  & 1.89 & 0.61\\
$\Delta$Debt in current liab (ew) & 100 & 4.9 & 0.40 &  & 0.03 & 0.32\\
\bottomrule
\end{tabular}


\end{table}


%%
% EXHIBIT-NUMBER: Table IA.5
\begin{table}
\caption{Themes from Mining Accounting Ratios in 2010}\label{tab:dm-theme2010}
\small

Table reports the 20 accounting ratio numerator and stock weight (equal- or value-) combinations with the largest mean t-stats using returns in the years 1963-2010 (IS).  `ew' is equal-weight, `vw' is value-weight. Strategies are signed to have positive mean returns IS. `Pct Short' is the share of strategies that short the ratio.  `t-stat' and  `Mean Return' are averages across the 65 possible denominators. 
\vspace{2ex}

\centering

\setlength{\tabcolsep}{1.0ex}

\begin{tabular}{lrrrlrr}
\toprule
 & \multicolumn{3}{c}{1963-2010 (IS)} &  & \multicolumn{1}{c}{2011-2014} & \multicolumn{1}{c}{2011-2022} \\ \cmidrule{2-4} \cmidrule{6-7}
Numerator (Stock Weight) & \multicolumn{1}{c}{Pct} & \multirow{2}{*}{t-stat} & \multicolumn{1}{c}{Mean} &  & \multicolumn{2}{c}{Mean Return} \\  
 & \multicolumn{1}{c}{Short} &  & \multicolumn{1}{c}{Return} &  & \multicolumn{2}{c}{OOS / IS} \\ 
\midrule
$\Delta$Long-term debt (ew) & 100 & 6.5 & 0.54 &  & 0.64 & 0.24\\
$\Delta$Inventories (ew) & 100 & 6.5 & 0.65 &  & 0.47 & 0.46\\
$\Delta$Liabilities (ew) & 100 & 6.4 & 0.68 &  & 0.52 & 0.14\\
$\Delta$Common stock (ew) & 100 & 6.3 & 0.60 &  & 0.24 & 0.40\\
$\Delta$Interest expense (ew) & 100 & 6.3 & 0.57 &  & 0.61 & 0.51\\
$\Delta$PPE net (ew) & 100 & 6.1 & 0.72 &  & 0.58 & 0.37\\
$\Delta$Current liabilities (ew) & 100 & 5.8 & 0.54 &  & 0.17 & 0.24\\
$\Delta$Debt convertible (ew) & 100 & 5.7 & 0.60 &  & 0.80 & 0.60\\
Merger sales contrib (ew) & 100 & 5.5 & 0.47 &  & 0.16 & 0.52\\
$\Delta$Assets (ew) & 100 & 5.5 & 0.81 &  & 0.40 & 0.35\\
$\Delta$Invested capital (ew) & 100 & 5.5 & 0.78 &  & 0.43 & 0.47\\
$\Delta$Intangible assets (ew) & 100 & 5.4 & 0.50 &  & 0.30 & 0.23\\
$\Delta$PPE gross (ew) & 100 & 5.3 & 0.65 &  & 0.52 & 0.45\\
$\Delta$Convertible debt and stock (ew) & 100 & 5.0 & 0.47 &  & 1.05 & 0.89\\
Retained earnings restatement (ew) & 100 & 5.0 & 0.51 &  & 0.53 & 0.19\\
$\Delta$Num employees (ew) & 100 & 4.9 & 0.54 &  & 0.25 & 0.53\\
$\Delta$Capital surplus (ew) & 100 & 4.8 & 0.65 &  & 0.54 & 0.97\\
$\Delta$Debt subordinated convertible (ew) & 100 & 4.8 & 0.63 &  & 0.34 & 0.86\\
$\Delta$Debt in current liab (ew) & 100 & 4.8 & 0.35 &  & 0.32 & 0.25\\
$\Delta$Capital expenditure (ew) & 100 & 4.7 & 0.47 &  & 0.36 & 0.89\\
\bottomrule
\end{tabular}


\end{table}


% =====================================================
\clearpage
\newpage
\section{Why Do Published Predictors Decay?}\label{sec:ia-decay}

Table \ref{tab:samp-split} illustrates two methods for documenting peer-reviewed predictability decay:
\begin{enumerate}
  \item Split at the end of the publication's sample period, following \citet{mclean2016does}
  \item Split in 2004 when high-speed internet became widely available, consistent with \citet{Chordia2014Have} and \citet{chen2022zeroing}
\end{enumerate}
Both approaches yield similar empirical results: a mean split date around 2000, a decay of about 50\%, with 85\% of predictors showing reduced effectiveness after the split.

%%
% EXHIBIT-NUMBER: Table IA.7
\begin{table}[!h]
  \caption{Why Do Peer-Reviewed Returns Decay?}\label{tab:samp-split}
  \small
  
  Table compares splitting samples using various methods: (1) the end of the original sample period, (2) when high speed internet became widely available, and (3) by minimizing the mean squared residual a la \citet{bai1998estimating}. Each method leads to a similar average break date, magnitude of decay, and frequency of decay.  It is unclear which sample split best explains why peer-reviewed predictability decays.
  \vspace{2ex}
  
  \centering
  \setlength{\tabcolsep}{1.0ex}
  
\begin{tabular}{lcccc}
\toprule
\multicolumn{1}{c}{Event} & \multicolumn{1}{c}{Mean} & \multicolumn{2}{c}{Return (bps p.m.)} & \multicolumn{1}{c}{\% of Signals} \\  & \multicolumn{1}{c}{Date} & \multicolumn{1}{c}{Before} & \multicolumn{1}{c}{After} & \multicolumn{1}{c}{w/ Decay} \\ 
\midrule
1. Paper's Sample Ends & Feb 2000 & 72 & 37 & 85\\
2. High Speed Internet & Dec 2004 & 71 & 31 & 88\\
3. Data-Driven Break & Mar 2001 & 80 & 25 & 82\\
\bottomrule
\end{tabular}

\end{table}

Which split best explains why peer-reviewed predictability decays? To examine this, we compute data driven breaks for each predictor by minimizing the mean squared residual (as in \citet{bai1998estimating}). We then compare the data-driven breaks with the breaks specified by the two methods above.

Figure \ref{fig:break-vs-sampend} shows the result. The scatter shows at best a mild relationship between the data-driven breaks and the papers' sample ends. Similarly, there is some clustering around the 2004 break, but the evidence is far from definitive. 

%%
% EXHIBIT-NUMBER: Figure IA.3
\begin{figure}[!h]
  \caption{Data-Driven Breaks vs Paper Sample Ends}
  \label{fig:break-vs-sampend}
  Each marker is one published predictor. Data-driven breaks split the predictor's sample into two periods to minimize the mean squared residual (as in \citet{bai1998estimating}).
  \vspace{0.15in}
  
  \centering
  \includegraphics[width=0.55\textwidth]{\exhibitspath/break_vs_sampend.pdf}
\end{figure}

This result is natural given the noise in long-short returns. The typical monthly volatility is 350 bps, implying the standard error of a 60-month mean is 45 bps, making it impossible to tell if a predictor decays  in a particular 5-year period.

Thus, it is difficult to determine the fundamental cause of the relative decay of peer-reviewed and data-mined predictors. This observation leads to our focus on making inferences about post-sample performance. 

\clearpage{}
\newpage{}
